% Lab 4 — compact version (full docs: docs/web/docs/labs/lab4/)
\section{Lab 4: Cryptography and Steganography CTF --- CipherStrike}

Lab~4 departs from the previous network-centric scenarios and challenges students
in the domain of applied cryptography and digital steganography. The fictional
premise is a forensic investigation at \textit{NovaCorp}: encrypted and steganographically
concealed communications have been intercepted following a corporate breach, and
the student acts as an incident response analyst tasked with recovering hidden data
and reconstructing the attacker's communication channel.

The lab is structured as a three-stage Capture-the-Flag (CTF): ServerA hosts five
classical and modern cryptography challenges; ServerB hosts five steganography
challenges; access to ServerC --- which hosts four advanced challenges combining
both disciplines --- is gated on completing key challenges from the first two modules.
The gate mechanism derives the ServerC SSH password from the MD5 hash of the
concatenation of two key flags:
\texttt{md5(A5\_flag + B5\_flag)[:12]}.
This design enforces a genuine learning dependency between stages rather than a
purely linear challenge sequence.
A quick teacher reference --- credential table, flag inventory, and startup procedure --- is compiled in Appendix~J;
complete challenge specifications and deployment guide are at \path{docs/web/docs/labs/lab4/}.


\begin{figure}[H]
\centering
\includegraphics[width=0.95\textwidth]{lab4-topology.png}
\caption{Lab~4 network topology}
\label{fig:lab4-topology}
\end{figure}

\subsubsection*{Network Addressing}

\begin{table}[H]
\centering
\resizebox{\textwidth}{!}{\begin{tabular}{@{}lllll@{}}
\toprule
\textbf{Node} & \textbf{Role} & \textbf{Interface} & \textbf{IP} & \textbf{Network} \\
\midrule
pfSense-LAB4 & Firewall/NAT   & WAN  & 192.168.0.18 (DHCP) & 192.168.0.0/24 (Homelab) \\
pfSense-LAB4 & Firewall       & LAN  & 192.168.1.1         & 192.168.1.0/24 \\
VyOS-LAB4    & Core router    & eth0 & 192.168.1.2         & 192.168.1.0/24 \\
VyOS-LAB4    & Zone router    & eth1 & 10.10.1.1           & 10.10.1.0/24 (Zone-A) \\
VyOS-LAB4    & Zone router    & eth2 & 10.10.2.1           & 10.10.2.0/24 (Zone-B) \\
VyOS-LAB4    & Zone router    & eth3 & 10.10.3.1           & 10.10.3.0/24 (Zone-C) \\
VyOS-LAB4    & Zone router    & eth4 & 10.10.4.1           & 10.10.4.0/24 (Zone-Defense) \\
ServerA      & Crypto module  & eth0 & 10.10.1.10          & 10.10.1.0/24 \\
ServerB      & Stego module   & eth0 & 10.10.2.10          & 10.10.2.0/24 \\
ServerC      & Advanced mix   & eth0 & 10.10.3.10          & 10.10.3.0/24 \\
Defender     & SOC monitor    & eth0 & 10.10.4.10          & 10.10.4.0/24 \\
\bottomrule
\end{tabular}}
\caption{Lab~4 IP addressing scheme}
\label{tab:lab4-ip}
\end{table}

\subsubsection*{CTF Challenge Progression}

\begin{table}[H]
\centering
\small
\begin{tabularx}{\textwidth}{lllX}
\toprule
\textbf{Server} & \textbf{Module} & \textbf{Challenges} & \textbf{Key technique categories} \\
\midrule
ServerA & Cryptography  & A1--A5 & ROT13, AES-ECB oracle, Vigen\`{e}re/Kasiski, XOR crib-dragging, RSA $e=3$ \\
ServerB & Steganography & B1--B5 & EXIF metadata, PNG chunk appending, LSB (RGB + alpha), steghide \\
ServerC & Advanced Mix  & C1--C4 & Multilayer stego, onion encoding, SHA-1 length extension, ECDSA nonce reuse \\
\bottomrule
\end{tabularx}
\caption{Lab~4 CTF module overview}
\label{tab:lab4-modules}
\end{table}

\subsection{Problems Encountered and Solutions}

\subsubsection{gmpy2 and pycryptodome Installation on Ubuntu 24.04}

Installing \texttt{gmpy2}~2.1.5 on Ubuntu Server~24.04 required the
\texttt{libgmp-dev}, \texttt{libmpfr-dev}, and \texttt{libmpc-dev} development
headers, which are not part of the default minimal image. The installation failed
silently when these were absent, producing a wheel build error rather than an
informative dependency message. The pre-staging script was updated to install the
native headers via \texttt{apt} before invoking \texttt{pip}, and the
\texttt{pycryptodome}~3.23.0 wheel was pinned to avoid upstream API changes
affecting the AES-ECB oracle challenge.

\subsubsection{Defender Node IP Address Persistence}

The Defender node is an Ubuntu Desktop instance used for monitoring. Its static
IP address (10.10.4.10) was not retained across EVE-NG lab restarts because the
NetworkManager connection profile was created interactively and not saved as a
persistent \texttt{.nmconnection} file. The issue was resolved by exporting the
active connection profile to \path{/etc/NetworkManager/system-connections/} and
setting the file permissions to \texttt{600}, after which the address survived
all subsequent restarts.

\subsubsection{hlextend Installation Dependency for ServerC}

The \texttt{hlextend} library required for the SHA-1 length extension challenge
(C3) is not available on PyPI and must be installed directly from its GitHub
repository. The pre-staging script therefore clones the repository and installs
it via \texttt{pip install -e}, which introduces a dependency on network access
during image preparation. To make image re-staging reproducible without internet
access on the EVE-NG host, the library source was vendored into the project
repository at \path{src/eve-ng/configs/nodes/Lab4/serverC/hlextend/} and the
deployment script installs from that local copy.
